\par\Needspace{4\baselineskip}
\begin{center}
\SourceLeaf{027c}\pdfbookmark[2]{10.}{article-010}10.
\end{center}
Die Gleichungen in Art. 8 enthalten alles, was zur Berechnung der exzentrischen und
der mittleren Anomalie aus der wahren oder der exzentrischen und der wahren Anomalie
aus der mittleren erforderlich ist. Um die exzentrische Anomalie aus der wahren zu
bestimmen, verwendet man gewöhnlich Formel VII. Meist eignet sich dazu jedoch Gl. X
besser, besonders wenn die Exzentrizität nicht zu groß ist. In diesem Fall lässt sich
$E$ nach X genauer berechnen als nach VII. Bei Anwendung von Gl. X erhält man außerdem
den in XII benötigten Logarithmus des Sinus von $E$ unmittelbar aus Gl. VIII. Bei
Anwendung von VII müsste man ihn den Tafeln entnehmen. Entnimmt man diesen Logarithmus
also auch beim erstgenannten Verfahren den Tafeln, so erhält man damit zugleich eine
Bestätigung dafür, dass die Rechnung richtig durchgeführt wurde. Solche Prüfungen
und Bestätigungen einer Rechnung sind stets von großem Wert. Wir werden daher bei
allen in diesem Werk dargelegten Verfahren, wo es sich ohne Weiteres einrichten lässt,
beständig für solche Kontrollen sorgen.
--- Zur weiteren Veranschaulichung fügen wir ein vollständig durchgerechnetes Beispiel bei.

\Needspace{22\baselineskip}
Gegeben seien $v = 310^0\,55^{\prime}\,29^{\prime\prime}{,}64$, $\varphi = 14^0\,12^{\prime}\,1^{\prime\prime}{,}87$ und $\log r = 0{,}3307640$. Gesucht sind $p$, $a$, $E$ und $M$.

\SourceLeaf{028}
\begin{center}
\begin{tabular}{@{}ll@{}}
\makebox[16em][l]{$\log\sin\varphi$\hfill .\,.\,.\,.\,.\,.\,.\,.\,.\,.} & $9{,}389\,7262$\\[.3em]
\makebox[16em][l]{$\log\cos v$\hfill .\,.\,.\,.\,.\,.\,.\,.\,.\,.} & $9{,}816\,2872$\\[.3em]
\hline
 & $9{,}206\,0134$\\[.3em]
\multicolumn{2}{r}{daraus $e\cos v=0{,}1606991$}\\[.3em]
\hline
\makebox[16em][l]{$\log(1+e\cos v)$\hfill .\,.\,.\,.\,.\,.} & $0{,}064\,7197$\\[.3em]
\makebox[16em][l]{$\log r$\hfill .\,.\,.\,.\,.\,.\,.\,.\,.\,.\,.\,.\,.} & $0{,}330\,7640$\\[.3em]
\hline
\makebox[16em][l]{$\log p$\hfill .\,.\,.\,.\,.\,.\,.\,.\,.\,.\,.\,.\,.} & $0{,}395\,4837$\\[.3em]
\makebox[16em][l]{$\log\cos\varphi^2$\hfill .\,.\,.\,.\,.\,.\,.\,.\,.} & $9{,}973\,0448$\\[.3em]
\hline
\makebox[16em][l]{$\log a$\hfill .\,.\,.\,.\,.\,.\,.\,.\,.\,.\,.\,.\,.} & $0{,}422\,4389$\\[.3em]
\hline
\makebox[16em][l]{$\log\sin v$\hfill .\,.\,.\,.\,.\,.\,.\,.\,.\,.} & $9{,}878\,2740_n$\footnotemark\\[.3em]
\makebox[16em][l]{$\log\BareRoot{\frac{p}{r}}$\hfill .\,.\,.\,.\,.\,.\,.\,.\,.\,.} & $0{,}032\,3598.5$\\[.3em]
\hline
 & $9{,}845\,9141.5_n$\\[.3em]
\makebox[16em][l]{$\log\sin\frac{1}{2}\varphi$\hfill .\,.\,.\,.\,.\,.\,.\,.\,.} & $9{,}092\,0395$\\[.3em]
\hline
\noalign{\vskip2pt}
\makebox[16em][l]{$\log\sin\frac{1}{2}(v-E)$\hfill .\,.\,.\,.\,.} & $8{,}937\,9536.5_n$\\[.3em]
\end{tabular}
\end{center}
\footnotetext{Der an einen Logarithmus angefügte Buchstabe $n$ zeigt an, dass die zugehörige Zahl negativ ist.}

\noindent Daraus folgt
\[
\begin{gathered}
\frac{1}{2}(v-E)=-4^0\,58^{\prime}\,22^{\prime\prime}{,}94;\quad v-E=-9^0\,56^{\prime}\,45^{\prime\prime}{,}88;\\
E=320^0\,52^{\prime}\,15^{\prime\prime}{,}52.
\end{gathered}
\]
\noindent Ferner ergibt sich
\begin{center}
\begin{tabular}{@{}ll@{}}
\makebox[16em][l]{$\log e$\hfill .\,.\,.\,.\,.\,.\,.\,.\,.} & $9{,}389\,7262$\\[.3em]
\makebox[16em][l]{$\log206264{,}8$\hfill .\,.\,.} & $5{,}314\,4251$\\[.3em]
\hline
\makebox[16em][l]{$\log e$ in Bogensek:\hfill .\,.\,.} & $4{,}704\,1513$\\[.3em]
\makebox[16em][l]{$\log\sin E$\hfill .\,.\,.\,.\,.\,.} & $9{,}800\,0767_n$\\[.3em]
\hline
 & $4{,}504\,2280_n$\\[.3em]
\end{tabular}
\end{center}

\begin{center}
Berechnung von $\log\sin E$ nach Formel VIII.
\end{center}
\begin{center}
\begin{tabular}{@{}ll@{}}
\makebox[16em][l]{$\log\frac{r}{p}\sin v$\hfill .\,.\,.\,.\,.} & $9{,}813\,5543_n$\\[.3em]
\makebox[16em][l]{$\log\cos\varphi$\hfill .\,.\,.\,.\,.\,.} & $9{,}986\,5224$\\[.3em]
\hline
\makebox[16em][l]{$\log\sin E$\hfill .\,.\,.\,.\,.\,.} & $9{,}800\,0767_n$\\[.3em]
\hline
\end{tabular}
\end{center}

\noindent Daraus folgt
\[
\begin{gathered}
e\sin E\ \mathrm{in\ Bogensekunden}=-31932^{\prime\prime}{,}14=-8^0\,52^{\prime}\,12^{\prime\prime}{,}14\\
\mathrm{und}\quad M=329^0\,44^{\prime}\,27^{\prime\prime}{,}66.
\end{gathered}
\]
\Needspace{12\baselineskip}
\noindent Mit Formel VII würde sich die Berechnung von $E$ wie folgt gestalten.
\begin{center}
\begin{tabular}{@{}ll@{}}
\multicolumn{2}{l}{$\frac{1}{2}v=155^0\,27^{\prime}\,44^{\prime\prime}{,}82$}\\[.3em]
\makebox[16em][l]{$\log\operatorname{tang}\frac{1}{2}v$\hfill .\,.\,.\,.\,.\,.\,.\,.} & $9{,}659\,4579_n$\\[.3em]
\multicolumn{2}{l}{$45^0-\frac{1}{2}\varphi=37^0\,53^{\prime}\,59^{\prime\prime}{,}065$}\\[.3em]
\makebox[16em][l]{$\log\operatorname{tang}(45^0-\frac{1}{2}\varphi)$\hfill .\,.\,.} & $9{,}891\,2427$\\[.3em]
\hline
\noalign{\vskip2pt}
\makebox[16em][l]{$\log\operatorname{tang}\frac{1}{2}E$\hfill .\,.\,.\,.\,.\,.\,.} & $9{,}550\,7006_n$\\[.3em]
\end{tabular}
\end{center}
Daraus erhält man $\frac{1}{2}E=160^0\,26^{\prime}\,7^{\prime\prime}{,}76$ und $E=320^0\,52^{\prime}\,15^{\prime\prime}{,}52$ wie oben.
